\section{Introdução}
\label{sec:introducao}

A qualidade interna de \emph{software} descreve propriedades estruturais do código-fonte que podem ser observadas sem a execução do programa, tais como modularidade, coesão, acoplamento e facilidade de manutenção \cite{iso25010,pressman:21}. Em sistemas orientados a objetos, essas propriedades têm sido operacionalizadas pela suíte de métricas proposta por Chidamber e Kemerer, que quantifica acoplamento, complexidade, herança e coesão em nível de classe \cite{chidamber:94}. No ecossistema Java, o \emph{framework} Spring Boot consolidou convenções de injeção de dependências, organização em camadas e autoconfiguração, tornando-se um recorte representativo para estudos empíricos sobre estrutura de código em repositórios públicos \cite{sun:22}. A combinação entre métricas clássicas e a ampla disponibilidade de projetos Spring Boot no GitHub permite examinar, com evidências extraídas do próprio código, como a comunidade de desenvolvimento organiza suas aplicações.

Apesar da consolidação dessas métricas e da popularidade do Spring Boot, a maior parte das investigações empíricas sobre qualidade estrutural em Java adota recortes transversais, isto é, analisa um conjunto de sistemas em um único instante, sem contrastar gerações distintas de projetos ao longo do tempo \cite{child:19,elemam:01}. Paralelamente, a partir de 2022, o desenvolvimento passou a incorporar de forma intensiva assistentes de inteligência artificial generativa, como GitHub Copilot, ChatGPT, Claude Code e Cursor, o que alterou o ambiente de produção de código \cite{peng:23,yetistiren:23}. \textbf{O problema investigado neste trabalho é a ausência de evidências empíricas longitudinais sobre como a qualidade estrutural de projetos Spring Boot hospedados no GitHub evolui ao longo de diferentes épocas históricas e se a Era da IA assistida introduziu alterações significativas nessas métricas.}

A relevância de tratar esse problema decorre de três fatores. Primeiro, a qualidade interna está associada ao esforço de manutenção e à viabilidade de longo prazo dos sistemas \cite{pressman:21,iso25010}; portanto, desconhecer a trajetória estrutural de um ecossistema tão difundido quanto o Spring Boot limita decisões de prática e de pesquisa. Segundo, a mineração de repositórios de \emph{software} no GitHub oferece base empírica em larga escala, desde que critérios de seleção, exclusão de \emph{forks} e maturidade temporal sejam observados \cite{hassan:08,kalliamvakou:16}. Terceiro, evidências recentes mostram que assistentes de IA reduzem o tempo de conclusão de tarefas e alteram a qualidade do código gerado em experimentos controlados \cite{peng:23,dakhel:23,liu:23}, mas ainda não esclarecem se esses efeitos se manifestam na estrutura de projetos reais publicados na plataforma. Sem uma análise por coortes, permanece indefinido se a Era da IA coincide com estabilidade, melhoria ou degradação das métricas estruturais.

O \textbf{objetivo geral deste trabalho é analisar longitudinalmente a evolução da qualidade estrutural de projetos Spring Boot \emph{open source} hospedados no GitHub, organizados em quatro coortes temporais, investigando o impacto da Era da IA assistida sobre métricas orientadas a objetos.} Os objetivos específicos são: (i)~caracterizar a evolução do acoplamento entre classes ao longo das quatro épocas; (ii)~caracterizar a evolução da complexidade e do tamanho de código; (iii)~caracterizar a evolução da herança e da coesão interna das classes; e (iv)~comparar estatisticamente o período imediatamente anterior à IA generativa (2018--2021) com a Era da IA assistida (2022--atual). A operacionalização desses objetivos no paradigma \emph{Goal-Question-Metric} é apresentada na Seção~\ref{sec:metodologia}.

Espera-se obter, como resultados, um retrato empírico das distribuições de acoplamento, complexidade, herança, coesão e tamanho em cada coorte, com testes não paramétricos que indiquem se as diferenças entre épocas são estatisticamente significativas, em especial no contraste entre o período pré-IA e a Era da IA assistida. Também se espera consolidar um pacote de replicação contendo \emph{scripts} de coleta, extração de métricas, notebooks de análise e artefatos de visualização, de modo a tornar o protocolo reproduzível no GitHub.

Este artigo está organizado da seguinte forma: a Seção~\ref{sec:fund} apresenta a fundamentação teórica; a Seção~\ref{sec:trabalhos} discute os trabalhos relacionados; a Seção~\ref{sec:metodologia} descreve o tipo de pesquisa, o ambiente, a amostragem, as métricas, os instrumentos, os experimentos já realizados e o cronograma do TCC~II; e o Apêndice~\ref{ap:visualizacoes} documenta o planejamento das visualizações longitudinais.
